% ЗГЕНЕРОВАНО scripts/build_format_v2_chapter.py — не редагувати вручну.
% Джерело: books/ashtavakra-gita/poetychnyi/rozdil-12.md

\poemverse{12.1}{Спершу нездатний до тілесних справ, \\ потім нездатний до поширення мови, \\ далі нездатний до думання — \\ тому саме в цьому я стою.}
\poemverse{12.2}{Через відсутність прихильності \\ до звуку й подібного \\ та через невидимість його для самого себе — \\ серце, для якого розсіяність і зосередженість одне: \\ саме в цьому я стою.}
\poemverse{12.3}{Коли розсіяність походить \\ від прагнення самадгі та подібного, \\ то й звичайна діяльність — задля самадгі. \\ Так побачивши це правило, \\ саме в цьому я стою.}
\poemverse{12.4}{Через відсутність того, що відкидати й приймати, \\ а отже — через відсутність радості й смутку, \\ нині, о Брахмане, \\ саме в цьому я стою.}
\poemverse{12.5}{Медитація зі стадією життя і без неї, \\ прийняття й відкидання, \\ зроблені умом, — \\ усе це домислювання. \\ Побачивши це в собі, \\ через саме це я так стою.}
\poemverse{12.6}{Виконання обрядів через незнання \\ і припинення їх — однаково. \\ Зрозумівши правильно цю суть, \\ я так стою.}
\poemverse{12.7}{Незбагненне, навіть коли його осмислюють, \\ набирає форми думки. \\ Тому, відкинувши те плекання, \\ я так стою.}
\poemverse{12.8}{У кого це так зроблено — \\ той стає тим, чию мету здійснено; \\ у кого така природа — \\ той стає тим, чию мету здійснено.}
